\documentclass[11pt, titlepage, letterpaper]{article}
\usepackage{style}

\begin{document}
\title{\LARGE{Comparison of Erasure-Aware Error Correction Decoding \\ for Fault Tolerant Quantum Computing \\}

    % \LARGE{Erasure is All You Need}

    \large{University of California, Berkeley\\C191A  Introduction to Quantum Computing I}\\
    Final Project. Fall 2025\\}

\author{Leigh Zaman\\ \texttt{leighzaman@berkeley.edu}
    \and Arin Sen \\ \texttt{arin.sen@berkeley.edu}
    \and Ilyes Tata \\ \texttt{ilyes.tata@berkeley.edu}
}
\date{Last revised:\\ \today}
\maketitle
\tableofcontents


\section{Introduction}

Motivation \& research question

Sub-questions \& hypotheses

\section{Background \& Related Work}

\subsection{Noise Models}
\paragraphsc{Pauli Noise.}
The simplest model used to evaluate Quantum Error Correction (QEC) assumes all physical qubit errors can be represented by uniformly random Pauli operators applied during circuit simulation. In our simulations we implement a \gls{depolarizing error channel} to represent \gls{Pauli noise}, which has the following effect on some quantum state $\rho$ with $p$ probability of becoming a maximally-mixed state and $(1-p)$ probability of remaining unechanged:

\begin{align}
    D(\rho, p) = (1-p)\rho + \dfrac{p}{3}(X\rho X + Y\rho Y + Z\rho Z) \quad \text{[Depolarizing channel]}
\end{align}

Unfortunately, error models based only on \gls{Pauli noise} are not robust enough to describe the real-world noise that must be corrected to realize fault-tolerant quantum computation. [CITATION NEEDED]

\paragraphsc{Leakage}
A \gls{leakage} ccurs when a qubit ``leaks'' out of the computational subspace (described by a two-dimensional Hilbert space) into a state that can only be described with a higher-order Hilbert space. If leakage goes undetected, the error compounds over time as it spreads through a multi-qubit interactions. Leakage-reduction units (LRUs) were proposed as a way to convert leakage errors into Pauli-type errors, based on a one-bit teleportation protocol \cite{suchara_leakage_2014}.

All species of qubit experience leakage errors, but it's important to note the underlying \emph{physical} processes look and behave very differently despite sharing a common mathematical description.

Leakage-detection units (LDUs) were developed to address leakage errors in superconducting qubits, by periodically coupling a data qubit to an acilla that flags when the data qubit leaves the computational subspace. The flagged data qubit can be reset/replaced during subsequent error-correction cycles, effectively treating the flagged qubit as an erasure error.

\paragraphsc{Correlated Errors.} TODO
\cite{wilen_correlated_2021}

\subsection{Erasure Qubit}

\paragraph{Transmons}: the dominant noise source is a form of detectable  Leakage-detection units (LDUs) were developed to address leakage errors in superconducting qubits, by periodically coupling a data qubit to an acilla that flags when the data qubit leaves the computational subspace. The flagged data qubit can be reset/replaced during subsequent error-correction cycles, effectively treating the flagged qubit as an erasure error.

A bias towards erasure-type errors occurs natively in some types of qubits (e.g. photons) or can be engineered by converting leakage errors to erasures (e.g. with LDUs in superconducting qubits).

The multileveled nature of neutral atom platforms leads to a complex error model. Neutral atoms are excited into ``Rydberg leakage'' describes an error that primarily occurs in multi-qubit gates. For example in two-qubit gates, a quantum mechanical phenomenon called the Rydberg blockade is exploited when exciting one atom to a high-energy Rydberg state, where at least one valence electron . The excited atom  when two or more atoms are excited into a high-energy Rydberg state typically when the biased error model, where the dominant type of error uccuring in two-qubit interactions is a leakage outside the computational subspace.

\subsection{Advantages of Erasure Errors}

\begin{enumerate}
    \item Can theoretically correct more erasure errors than unknown Pauli errors.
          We can correct $d-1$ erasure errors compared to $\dfrac{d-1}{2}$ Pauli errors. As a result, we expect logical error rate $P_L$ on the order of $p^{d}$ where $p$ is the physical qubit error rate, compared to $p^{d/2}$ for unknown Pauli errors. \cite{gu_2023}
    \item Higher thresholds: surface code 50\% erasure noise vs. 19\% depolarizing noise.
    \item Third item
\end{enumerate}


Figure: d=5 surface code with some unknown Pauli error spanning 3 qubits, with syndrome squares highlighted. If we apply the minimum-weight correction, this causes a logical error.

Figure: d=5 surface code with erasure errors, exact locations known.



\section{Methods}

\section{Metrics}

\section{Results}

\section{Discussion}

\section{Conclusion}

\section{Acknowledgements}


\newpage
% cite all references
% \nocite{*}
\printbibliography[title={References},heading=bibnumbered]

\newpage
\appendix

\section{Appendix: Syndrome of Loss Errors}

Describe how a loss is heralded during state readout. Each potential loss location is described by a distinct error probability, described by a unique configuration of the error syndroms. [EXAMPLE OF ERROR SYNDROME]

\section{Appendix: Additional Figures and Tables}

\newpage
\printglossaries

\end{document}